\documentclass[a4paper]{arrowhead}

\usepackage[yyyymmdd]{datetime}
\usepackage{etoolbox}
\usepackage[utf8]{inputenc}
\usepackage[T1]{fontenc}
\usepackage{multirow}
\usepackage{longtable}
\usepackage{tikz}
\usetikzlibrary{positioning,shapes.multipart,arrows.meta}

\renewcommand{\dateseparator}{-}
\setlength{\emergencystretch}{3em}
\lstset{language=HTTP,breakatwhitespace=false,columns=fullflexible}

\newcommand{\fparam}[1]{\textit{\textcolor{ArrowheadBlue}{#1}}}

%% Special references
\newcommand{\mref}[2][]{{\textcolor{ArrowheadPurple}{\hyperref[sec:model:#2]{\ifstrempty{#1}{#2}{#1}}}}}
\newcommand{\oref}[1]{{\textcolor{ArrowheadBlue}{\hyperref[sec:operations:#1]{#1}}}}
\newcommand{\pdef}[1]{{\textcolor{ArrowheadGrey}{#1 \label{sec:model:primitives:#1} \label{sec:model:primitives:#1s}}}}
\newcommand{\pref}[1]{{\textcolor{ArrowheadGrey}{\hyperref[sec:model:primitives:#1]{#1}}}}

% as template_idd.tex; the input and output arguments are passed already formatted (\mref or \pref)
\newrobustcmd\osubsection[5]{
  \addtocounter{subsection}{1}
  \addcontentsline{toc}{subsection}{\protect\numberline{\thesubsection}operation \textcolor{ArrowheadBlue}{#1}}
  \subsection*{
    \thesubsection\quad
    #2 \textcolor{ArrowheadPurple}{#3} \\
    \small
    \hspace*{0.075\textwidth}\begin{minipage}{0.1\textwidth}
      \vspace*{1mm}
      Operation: \\
      \notblank{#4}{Input: \\}{}
      \notblank{#5}{Output: \\}{}
    \end{minipage}
    \begin{minipage}{0.825\textwidth}
      \vspace*{1mm}
      \textcolor{ArrowheadBlue}{#1} \\
      \notblank{#4}{#4 \\}{}
      \notblank{#5}{#5 \\}{}
    \end{minipage}
  }
  \label{sec:operations:#1}
}
\newrobustcmd\msubsection[2]{
  \addtocounter{subsection}{1}
  \addcontentsline{toc}{subsection}{\protect\numberline{\thesubsection}#1 \textcolor{ArrowheadPurple}{#2}}
  \subsection*{\thesubsection\quad#1 \textcolor{ArrowheadPurple}{#2}}
  \label{sec:model:#2} \label{sec:model:#2s}
}
%%

\begin{document}
\sloppy

%% Arrowhead Document Properties
\ArrowheadTitle{translationBridge MQTT/TLS/JSON}
\ArrowheadServiceID{translationBridge}
\ArrowheadType{Interface Design Description}
\ArrowheadTypeShort{IDD}
\ArrowheadVersion{5.2.2}
\ArrowheadDate{2026-10-30}
\ArrowheadAuthor{Jerker Delsing}
\ArrowheadStatus{RELEASE CANDIDATE}
\ArrowheadContact{jerker.delsing@ltu.se}
\ArrowheadFooter{\href{www.arrowhead.eu}{www.arrowhead.eu}}
\ArrowheadSetup
%%

%% Front Page
\begin{center}
  \vspace*{1cm}
  \huge{\arrowtitle}

  \vspace*{0.2cm}
  \LARGE{\arrowtype}
  \vspace*{1cm}

  \Large{Service ID: \textit{"\arrowid"}}

  \vspace*{\fill}

  % Front Page Image

  \vspace*{1cm}
  \vspace*{\fill}

  % Front Page Abstract
  \begin{abstract}
    This document describes the MQTT protocol with TLS payload security and JSON payload encoding variant of the \texttt{translationBridge} service of the TranslationManager system.
  \end{abstract}

  \vspace*{1cm}
\end{center}

\newpage
%%

%% Table of Contents
\tableofcontents
\newpage
%%

\section{Overview}
\label{sec:overview}

This document describes the \texttt{translationBridge} service interface of the TranslationManager system over MQTT, which provides interface and data-model translation for a consumer: finding providers that become usable through a translation bridge, building the bridge and aborting it.

It's implemented using protocol, encoding as stated in below table:
\begin{table}[ht!]
  \centering
  \begin{tabular}{|l|l|l|}
    \hline
    \rowcolor{gray!33} Profile type & Type & Version \\ \hline
    Transfer protocol & MQTT & 3.1.1 \cite{mqtt311} \\ \hline
    Data encryption & TLS & 1.3 \cite{rfc8446}, 1.2 \cite{rfc5246} \\ \hline
    Encoding & JSON & RFC 8259 \cite{rfc8259} \\ \hline
    Compression & N/A & - \\ \hline
    Semantics & N/A & - \\ \hline
    Ontology & N/A & - \\ \hline
  \end{tabular}
  \caption{Communication and semantics details used for the \texttt{translationBridge} service interface}
  \label{tab:comunication_semantics_profile}
\end{table}

This document provides the Interface Design Description IDD to the \textit{translationBridge -- Service Description} document \cite{tm-translationBridge-sd}.
The operations, their semantics and the data models are the same as in the HTTP/TLS/JSON IDD \cite{tm-translationBridge-http-idd}; this document specifies the MQTT binding.

The MQTT interface is offered only when the system runs with \texttt{mqtt.api.enabled=true}; it is then registered with the interface template \texttt{generic\_\allowbreak{}mqtts} (\texttt{generic\_\allowbreak{}mqtt} without TLS), the broker address and port (\texttt{mqtt.\allowbreak{}broker.\allowbreak{}address}, \texttt{mqtt.\allowbreak{}broker.\allowbreak{}port}) and the base topic \texttt{arrowhead/\allowbreak{}translation/\allowbreak{}bridge/\allowbreak{}}. The standalone deployment sets \texttt{mqtt.\allowbreak{}api.\allowbreak{}enabled=\allowbreak{}false}, so the MQTT interfaces are not registered there.

\newpage

\section{Service Operations}
\label{sec:operations}

The provider subscribes to one topic per operation, the base topic followed by the operation name:

\noindent\begin{longtable}{| p{5cm} | p{10.4cm} |} \hline
\rowcolor{gray!33} Operation & Request topic \\ \hline \endhead
\texttt{discovery} & \texttt{arrowhead/\allowbreak{}translation/\allowbreak{}bridge/\allowbreak{}discovery} \\ \hline
\texttt{negotiation} & \texttt{arrowhead/\allowbreak{}translation/\allowbreak{}bridge/\allowbreak{}negotiation} \\ \hline
\texttt{abort} & \texttt{arrowhead/\allowbreak{}translation/\allowbreak{}bridge/\allowbreak{}abort} \\ \hline
\end{longtable}

A consumer publishes a request envelope (\mref{MqttRequestTemplate}) to the topic and receives the answer (\mref{MqttResponseTemplate}) on its own \texttt{responseTopic}. The \texttt{authentication} field carries the requester's identity according to the provider's authentication policy: \texttt{IDENTITY-TOKEN//<token>} (outsourced or internal policy), \texttt{SYSTEM//<name>} (declared policy) or the requester's X.509 certificate (certificate policy). The \texttt{status} of the response uses the HTTP status code numbers (200, 201, 204, 400, 401, 403, 404, 408, 423, 500, 503); on error the payload is an \mref{ErrorMessage}.

\begin{lstlisting}[caption={Request envelope.}]
{
  "traceId": "a1b2c3",
  "authentication": "IDENTITY-TOKEN//<token>",
  "responseTopic": "arrowhead/response/TemperatureDisplay/a1b2c3",
  "qosRequirement": 1,
  "params": {},
  "payload": {}
}
\end{lstlisting}

\begin{lstlisting}[caption={Response envelope.}]
{
  "status": 200,
  "traceId": "a1b2c3",
  "receiver": "TemperatureDisplay",
  "payload": {}
}
\end{lstlisting}

\osubsection{discovery}{PUBLISH}{\texttt{arrowhead/\allowbreak{}translation/\allowbreak{}bridge/\allowbreak{}discovery}}{\mref[Translation\allowbreak{}Discovery\allowbreak{}Request]{TranslationDiscoveryRequest}}{\mref[Translation\allowbreak{}Discovery\allowbreak{}Response]{TranslationDiscoveryResponse}}

Returns the candidates usable by the requester through translation, with a bridge identifier.

\begin{lstlisting}[caption={Payload of an \oref{discovery} request.}]
{
  "candidates": [
    {
      "instanceId": "TemperatureSensor|temperatureReading|1.0.0",
      "provider": {
        "name": "TemperatureSensor",
        "metadata": {
          "location": "hall-1"
        },
        "version": "1.0.0",
        "addresses": [
          {}
        ],
        "device": {
          "name": "Raspberry1",
          "metadata": {
            "location": "hall-1"
          },
          "addresses": [
            {}
          ],
          "createdAt": "2026-10-22T12:00:00Z",
          "updatedAt": "2026-10-22T12:00:00Z"
        },
        "createdAt": "2026-10-22T12:00:00Z",
        "updatedAt": "2026-10-22T12:00:00Z"
      },
      "serviceDefinition": {
        "name": "temperatureReading",
        "createdAt": "2026-10-22T12:00:00Z",
        "updatedAt": "2026-10-22T12:00:00Z"
      },
      "version": "1.0.0",
      "expiresAt": "2026-10-22T12:00:00Z",
      "metadata": {
        "unit": "celsius"
      },
      "interfaces": [
        {
          "templateName": "generic_https",
          "protocol": "https",
          "policy": "CERT_AUTH",
          "properties": {
            "accessAddresses": [
              "192.168.0.10"
            ],
            "accessPort": 8080,
            "basePath": "/temperature",
            "operations": {
              "read": {
                "method": "GET",
                "path": "/read"
              }
            }
          }
        }
      ],
      "createdAt": "2026-10-22T12:00:00Z",
      "updatedAt": "2026-10-22T12:00:00Z"
    }
  ],
  "operation": "FILTER",
  "interfaceTemplateNames": [
    "string"
  ],
  "inputDataModelId": "3fa85f64-5717-4562-b3fc-2c963f66afa6",
  "outputDataModelId": "3fa85f64-5717-4562-b3fc-2c963f66afa6"
}
\end{lstlisting}

\osubsection{negotiation}{PUBLISH}{\texttt{arrowhead/\allowbreak{}translation/\allowbreak{}bridge/\allowbreak{}negotiation}}{\mref[Translation\allowbreak{}Negotiation\allowbreak{}Request]{TranslationNegotiationRequest}}{\mref[Translation\allowbreak{}Negotiation\allowbreak{}Response]{TranslationNegotiationResponse}}

Builds the bridge to the chosen candidate and returns the bridge interface and token limits.

\begin{lstlisting}[caption={Payload of an \oref{negotiation} request.}]
{
  "bridgeId": "3fa85f64-5717-4562-b3fc-2c963f66afa6",
  "target": {
    "instanceId": "TemperatureSensor|temperatureReading|1.0.0",
    "provider": {
      "name": "TemperatureSensor",
      "metadata": {
        "location": "hall-1"
      },
      "version": "1.0.0",
      "addresses": [
        {
          "type": "IPV4",
          "address": "192.168.0.10"
        }
      ],
      "device": {
        "name": "Raspberry1",
        "metadata": {
          "location": "hall-1"
        },
        "addresses": [
          {}
        ],
        "createdAt": "2026-10-22T12:00:00Z",
        "updatedAt": "2026-10-22T12:00:00Z"
      },
      "createdAt": "2026-10-22T12:00:00Z",
      "updatedAt": "2026-10-22T12:00:00Z"
    },
    "serviceDefinition": {
      "name": "temperatureReading",
      "createdAt": "2026-10-22T12:00:00Z",
      "updatedAt": "2026-10-22T12:00:00Z"
    },
    "version": "1.0.0",
    "expiresAt": "2026-10-22T12:00:00Z",
    "metadata": {
      "unit": "celsius"
    },
    "interfaces": [
      {
        "templateName": "generic_https",
        "protocol": "https",
        "policy": "CERT_AUTH",
        "properties": {
          "accessAddresses": [
            "192.168.0.10"
          ],
          "accessPort": 8080,
          "basePath": "/temperature",
          "operations": {
            "read": {
              "method": "GET",
              "path": "/read"
            }
          }
        }
      }
    ],
    "createdAt": "2026-10-22T12:00:00Z",
    "updatedAt": "2026-10-22T12:00:00Z"
  },
  "operation": "FILTER",
  "interfaceTemplateName": "string",
  "inputDataModelId": "3fa85f64-5717-4562-b3fc-2c963f66afa6",
  "outputDataModelId": "3fa85f64-5717-4562-b3fc-2c963f66afa6"
}
\end{lstlisting}

\osubsection{abort}{PUBLISH}{\texttt{arrowhead/\allowbreak{}translation/\allowbreak{}bridge/\allowbreak{}abort}}{\pref{String}}{}

Aborts a bridge of the requester. Returns 200 if aborted, 204 if there was nothing to abort. The bridge identifier fills the \texttt{\{bridgeId\}} placeholder of the registered path.

\begin{lstlisting}[caption={Payload of an \oref{abort} request.}]
"3fa85f64-5717-4562-b3fc-2c963f66afa6"
\end{lstlisting}

\newpage

\section{Data Models}
\label{sec:model}

Here, all data objects that can be part of the service calls associated with this service are listed in alphabetic order.
Each \textit{struct} denotes a JSON \pref{Object} with the named fields; fields that are not needed may be omitted.

\msubsection{struct}{MqttRequestTemplate}

The request envelope.

\noindent\begin{tabularx}{\textwidth}{| p{4.25cm} | p{3.5cm} | X |} \hline
\rowcolor{gray!33} Object Field & Value Type & Description \\ \hline
"traceId" & \pref{String} & Identifier chosen by the requester, returned in the response. \\ \hline
"authentication" & \pref{String} & Identity of the requester (see Section \ref{sec:operations}). \\ \hline
"responseTopic" & \pref{String} & Topic on which the response is published. \\ \hline
"qosRequirement" & \pref{Number} & MQTT QoS of the response (0, 1 or 2). \\ \hline
"params" & \pref{Object}$<$\pref{String}$>$ & Optional parameters (the query parameters of the HTTP binding). \\ \hline
"payload" & \pref{Object} & The operation input. \\ \hline
\end{tabularx}

\msubsection{struct}{MqttResponseTemplate}

The response envelope.

\noindent\begin{tabularx}{\textwidth}{| p{4.25cm} | p{3.5cm} | X |} \hline
\rowcolor{gray!33} Object Field & Value Type & Description \\ \hline
"status" & \pref{Number} & Status code. \\ \hline
"traceId" & \pref{String} & The trace identifier of the request. \\ \hline
"receiver" & \pref{String} & Name of the requester. \\ \hline
"payload" & \pref{Object} & The operation output, or an \mref{ErrorMessage}. \\ \hline
\end{tabularx}

\msubsection{struct}{ErrorMessage}




\noindent\begin{tabularx}{\textwidth}{| >{\raggedright\arraybackslash}p{3.9cm} | >{\raggedright\arraybackslash\small}p{5.3cm} | X |} \hline
\rowcolor{gray!33} Object Field & Value Type & Description \\ \hline
"error\allowbreak{}Message" & \pref{String} & Error text, if the event is an error. \\ \hline
"error\allowbreak{}Code" & \pref{Number} &  \\ \hline
"exception\allowbreak{}Type" & \pref{ExceptionType} &  \\ \hline
"origin" & \pref{String} &  \\ \hline
\end{tabularx}



\msubsection{struct}{Address}

A network address with its type.


\noindent\begin{tabularx}{\textwidth}{| >{\raggedright\arraybackslash}p{3.9cm} | >{\raggedright\arraybackslash\small}p{5.3cm} | X |} \hline
\rowcolor{gray!33} Object Field & Value Type & Description \\ \hline
"type" & \pref{String} & Address type: \texttt{HOSTNAME}, \texttt{IPV4}, \texttt{IPV6} or \texttt{MAC}. \\ \hline
"address" & \pref{String} & One network address (host name, IPv4, IPv6 or MAC address). \\ \hline
\end{tabularx}



\msubsection{struct}{DeviceResponse}

A registered device.


\noindent\begin{tabularx}{\textwidth}{| >{\raggedright\arraybackslash}p{3.9cm} | >{\raggedright\arraybackslash\small}p{5.3cm} | X |} \hline
\rowcolor{gray!33} Object Field & Value Type & Description \\ \hline
"name" & \pref{String} & Device name. \\ \hline
"metadata" & \pref{Object} & Free-form metadata (JSON object). \\ \hline
"addresses" & \pref{Array}$<$\allowbreak{}\mref{Address}$>$ & Network addresses with their types. \\ \hline
"created\allowbreak{}At" & \pref{DateTime} & UTC time of creation. \\ \hline
"updated\allowbreak{}At" & \pref{DateTime} & UTC time of the last change. \\ \hline
\end{tabularx}



\msubsection{struct}{ServiceDefinitionResponse}

A service definition.


\noindent\begin{tabularx}{\textwidth}{| >{\raggedright\arraybackslash}p{3.9cm} | >{\raggedright\arraybackslash\small}p{5.3cm} | X |} \hline
\rowcolor{gray!33} Object Field & Value Type & Description \\ \hline
"name" & \pref{String} & Service definition name. \\ \hline
"created\allowbreak{}At" & \pref{DateTime} & UTC time of creation. \\ \hline
"updated\allowbreak{}At" & \pref{DateTime} & UTC time of the last change. \\ \hline
\end{tabularx}



\msubsection{struct}{ServiceInstanceInterfaceResponse}

One interface of a service instance.


\noindent\begin{tabularx}{\textwidth}{| >{\raggedright\arraybackslash}p{3.9cm} | >{\raggedright\arraybackslash\small}p{5.3cm} | X |} \hline
\rowcolor{gray!33} Object Field & Value Type & Description \\ \hline
"template\allowbreak{}Name" & \pref{String} & Name of the interface template, snake\_case (e.g. \texttt{generic\_\allowbreak{}https}). \\ \hline
"protocol" & \pref{String} & Protocol of the interface, e.g. \texttt{http}, \texttt{https}, \texttt{mqtt}, \texttt{mqtts}. \\ \hline
"policy" & \pref{String} & Security policy of the interface: \texttt{NONE}, \texttt{CERT\_\allowbreak{}AUTH} or one of the token policies. \\ \hline
"properties" & \pref{Object} & Interface properties defined by the template, e.g. \texttt{accessAddresses}, \texttt{accessPort}, \texttt{basePath}, \texttt{operations}. \\ \hline
\end{tabularx}



\msubsection{struct}{ServiceInstanceResponse}

A registered service instance.


\noindent\begin{tabularx}{\textwidth}{| >{\raggedright\arraybackslash}p{3.9cm} | >{\raggedright\arraybackslash\small}p{5.3cm} | X |} \hline
\rowcolor{gray!33} Object Field & Value Type & Description \\ \hline
"instance\allowbreak{}Id" & \pref{String} & Service instance identifier \texttt{\textless{}System\textgreater{}\textbar{}\allowbreak{}\textless{}service\allowbreak{}Definition\textgreater{}\textbar{}\allowbreak{}\textless{}version\textgreater{}}, or policy instance identifier. \\ \hline
"provider" & \mref{SystemResponse} & The provider system. \\ \hline
"service\allowbreak{}Definition" & \mref[Service\allowbreak{}Definition\allowbreak{}Response]{ServiceDefinitionResponse} & The service definition. \\ \hline
"version" & \pref{Version} & Version of the service instance. \\ \hline
"expires\allowbreak{}At" & \pref{DateTime} & UTC expiry time; no value means no expiry. \\ \hline
"metadata" & \pref{Object} & Free-form metadata (JSON object). \\ \hline
"interfaces" & \pref{Array}$<$\allowbreak{}\mref[Service\allowbreak{}Instance\allowbreak{}Interface\allowbreak{}Response]{ServiceInstanceInterfaceResponse}$>$ & Service interfaces offered by the instance. \\ \hline
"created\allowbreak{}At" & \pref{DateTime} & UTC time of creation. \\ \hline
"updated\allowbreak{}At" & \pref{DateTime} & UTC time of the last change. \\ \hline
\end{tabularx}



\msubsection{struct}{SystemResponse}

A registered system.


\noindent\begin{tabularx}{\textwidth}{| >{\raggedright\arraybackslash}p{3.9cm} | >{\raggedright\arraybackslash\small}p{5.3cm} | X |} \hline
\rowcolor{gray!33} Object Field & Value Type & Description \\ \hline
"name" & \pref{String} & System name. \\ \hline
"metadata" & \pref{Object} & Free-form metadata (JSON object). \\ \hline
"version" & \pref{Version} & Version of the system. \\ \hline
"addresses" & \pref{Array}$<$\allowbreak{}\mref{Address}$>$ & Network addresses with their types. \\ \hline
"device" & \mref{DeviceResponse} & The device hosting the system, if registered. \\ \hline
"created\allowbreak{}At" & \pref{DateTime} & UTC time of creation. \\ \hline
"updated\allowbreak{}At" & \pref{DateTime} & UTC time of the last change. \\ \hline
\end{tabularx}



\msubsection{struct}{TranslationBridgeCandidate}

A candidate reachable through translation.


\noindent\begin{tabularx}{\textwidth}{| >{\raggedright\arraybackslash}p{3.9cm} | >{\raggedright\arraybackslash\small}p{5.3cm} | X |} \hline
\rowcolor{gray!33} Object Field & Value Type & Description \\ \hline
"service\allowbreak{}Instance\allowbreak{}Id" & \pref{String} & Service instance identifier \texttt{\textless{}System\textgreater{}\textbar{}\allowbreak{}\textless{}service\allowbreak{}Definition\textgreater{}\textbar{}\allowbreak{}\textless{}version\textgreater{}}. \\ \hline
"interface\allowbreak{}Template\allowbreak{}Name" & \pref{String} & Interface template of the candidate the bridge would translate to. \\ \hline
\end{tabularx}



\msubsection{struct}{TranslationDiscoveryRequest}

Discovery of translation bridges for the requester.


\noindent\begin{tabularx}{\textwidth}{| >{\raggedright\arraybackslash}p{3.9cm} | >{\raggedright\arraybackslash\small}p{5.3cm} | X |} \hline
\rowcolor{gray!33} Object Field & Value Type & Description \\ \hline
"candidates" & \pref{Array}$<$\allowbreak{}\mref{ServiceInstanceResponse}$>$ & Candidate service instances (from service discovery). \\ \hline
"operation" & \pref{String} & Operation the consumer wants to call. \\ \hline
"interface\allowbreak{}Template\allowbreak{}Names" & \pref{Array}$<$\allowbreak{}\pref{String}$>$ & Interface templates the consumer can use. \\ \hline
"input\allowbreak{}Data\allowbreak{}Model\allowbreak{}Id" & \pref{String} & Data model of the operation input expected by the consumer. \\ \hline
"output\allowbreak{}Data\allowbreak{}Model\allowbreak{}Id" & \pref{String} & Data model of the operation output expected by the consumer. \\ \hline
\end{tabularx}



\msubsection{struct}{TranslationDiscoveryResponse}

Candidates reachable through a translation bridge.


\noindent\begin{tabularx}{\textwidth}{| >{\raggedright\arraybackslash}p{3.9cm} | >{\raggedright\arraybackslash\small}p{5.3cm} | X |} \hline
\rowcolor{gray!33} Object Field & Value Type & Description \\ \hline
"bridge\allowbreak{}Id" & \pref{String} & Identifier (UUID) of the translation bridge. \\ \hline
"candidates" & \pref{Array}$<$\allowbreak{}\mref[Translation\allowbreak{}Bridge\allowbreak{}Candidate]{TranslationBridgeCandidate}$>$ & Store entries to create. \\ \hline
\end{tabularx}



\msubsection{struct}{TranslationNegotiationRequest}

Builds the bridge to the chosen candidate.


\noindent\begin{tabularx}{\textwidth}{| >{\raggedright\arraybackslash}p{3.9cm} | >{\raggedright\arraybackslash\small}p{5.3cm} | X |} \hline
\rowcolor{gray!33} Object Field & Value Type & Description \\ \hline
"bridge\allowbreak{}Id" & \pref{String} & Identifier (UUID) of the translation bridge. \\ \hline
"target" & \mref{ServiceInstanceResponse} & The chosen candidate service instance. \\ \hline
"operation" & \pref{String} & QoS operation: \texttt{FILTER} or \texttt{SORT}. \\ \hline
"interface\allowbreak{}Template\allowbreak{}Name" & \pref{String} & Interface template the consumer will use. \\ \hline
"input\allowbreak{}Data\allowbreak{}Model\allowbreak{}Id" & \pref{String} & Data model of the operation input expected by the consumer. \\ \hline
"output\allowbreak{}Data\allowbreak{}Model\allowbreak{}Id" & \pref{String} & Data model of the operation output expected by the consumer. \\ \hline
\end{tabularx}



\msubsection{struct}{TranslationNegotiationResponse}

The established bridge.


\noindent\begin{tabularx}{\textwidth}{| >{\raggedright\arraybackslash}p{3.9cm} | >{\raggedright\arraybackslash\small}p{5.3cm} | X |} \hline
\rowcolor{gray!33} Object Field & Value Type & Description \\ \hline
"bridge\allowbreak{}Id" & \pref{String} & Identifier (UUID) of the translation bridge. \\ \hline
"bridge\allowbreak{}Interface" & \mref[Service\allowbreak{}Instance\allowbreak{}Interface\allowbreak{}Response]{ServiceInstanceInterfaceResponse} & Interface of the bridge the consumer calls instead of the provider. \\ \hline
"token\allowbreak{}Expires\allowbreak{}At" & \pref{DateTime} & Expiry of the bridge token. \\ \hline
"token\allowbreak{}Usage\allowbreak{}Limit" & \pref{Number} & Number of uses allowed with the bridge token. \\ \hline
\end{tabularx}



\subsection{Primitives}
\label{sec:model:primitives}

As all messages are encoded using the JSON format \cite{rfc8259}, the following primitive constructs, part of that standard, become available.

\noindent\begin{tabularx}{\textwidth}{| p{4.6cm} | X |} \hline
\rowcolor{gray!33} Type & Description \\ \hline
\pdef{Array} & A JSON Array \cite{rfc8259}; \pref{Array}$<$A$>$ has elements of type A. \\ \hline
\pdef{DateTime} & JSON String with an ISO 8601 UTC timestamp, e.g. \texttt{2026-10-22T12:00:00Z}. \\ \hline
\pdef{ExceptionType} & One of \texttt{ARROWHEAD}, \texttt{INVALID\_\allowbreak{}PARAMETER}, \texttt{AUTH}, \texttt{FORBIDDEN}, \texttt{DATA\_\allowbreak{}NOT\_\allowbreak{}FOUND}, \texttt{TIMEOUT}, \texttt{LOCKED}, \texttt{INTERNAL\_\allowbreak{}SERVER\_\allowbreak{}ERROR}, \texttt{EXTERNAL\_\allowbreak{}SERVER\_\allowbreak{}ERROR} (JSON String). \\ \hline
\pdef{Number} & JSON Number, an integer. \\ \hline
\pdef{Object} & JSON Object; \pref{Object}$<$A$>$ has string keys and values of type A. \\ \hline
\pdef{String} & JSON String. \\ \hline
\pdef{Version} & JSON String, semantic version \texttt{major.minor.patch}. \\ \hline
\end{tabularx}

\newpage

\bibliographystyle{IEEEtran}
\bibliography{bibliography}

\newpage

\section{Revision History}
\subsection{Amendments}

\noindent\begin{tabularx}{\textwidth}{| p{1cm} | p{3cm} | p{2cm} | X | p{4cm} |} \hline
\rowcolor{gray!33} No. & Date & Version & Subject of Amendments & Author \\ \hline

1 & 2026-10-22 & \arrowversion & Document created from the Arrowhead LaTeX templates for the 5.2.2 Java Spring implementation, including the local extensions of this workspace. & Jerker Delsing \\ \hline
2 & 2026-10-30 & \arrowversion & Checked against the primary Java source tree (Eclipse-Arrowhead-java): standalone deployment, configuration and security statements updated; missing primitive definitions added. & Claude Opus 5.5 (draft) \\ \hline
\end{tabularx}

\subsection{Quality Assurance}

\noindent\begin{tabularx}{\textwidth}{| p{1cm} | p{3cm} | p{2cm} | X |} \hline
\rowcolor{gray!33} No. & Date & Version & Approved by \\ \hline

1 & & \arrowversion  &  \\ \hline

\end{tabularx}

\end{document}
